ગણ $A$ પરનો \emph{સંબંધ} $R$ એ $A$ના ઘટકો વચ્ચે સંબંધ સ્થાપિત કરવાની રીત છે. $x$ અને $y$ વચ્ચે $R$ લાગુ પડતો હોય તો $Rxy$ લખીએ છીએ. ઔપચારિક રીતે $R$ એ $Rxy$ હોય એવી જોડો $\tuple{x,y}\in A^2$નો ગણ છે. નાનું હોવું, મોટું હોવું, સરખું હોવું, નિઃશેષ ભાગ આપવો, સમાન લંબાઈ હોવી, ઉપગણ હોવો અને સમાન કદ હોવું એ સંખ્યાઓ, શ્રેણીઓ કે ગણો પરના સંબંધોનાં અગત્યનાં ઉદાહરણો છે. સંબંધો ચિત્રરૂપે બતાવવા \emph{આલેખો} વપરાય છે. આલેખને \emph{શીર્ષો}ના ગણ અને તેના પરના દ્વિઘટકી \emph{ધાર} સંબંધ તરીકે પણ સમજી શકાય છે.

કેટલાક સંબંધોના સામાન્ય ગુણધર્મ તેમને ખાસ ઉપયોગી બનાવે છે. દરેક ઘટક પોતાની સાથે $R$-સંબંધમાં હોય તો $R$ \emph{સ્વવાચક} છે. કોઈ પણ $x,y$ માટે $Rxy$ હોય ત્યારે $Ryx$ પણ હોય તો તે \emph{સંમિત} છે. $Rxy$ અને $Ryz$ પરથી $Rxz$ મળે તો તે \emph{પરંપરિત} છે. આ ત્રણેય ગુણધર્મો ધરાવતા સંબંધો \emph{સામ્ય સંબંધો} છે. $Rxy$ અને $Ryx$ પરથી $x=y$ મળે તો સંબંધ \emph{વિસંમિત} છે. સ્વવાચક, વિસંમિત અને પરંપરિત સંબંધ \emph{આંશિક ક્રમ} છે. કોઈ પણ $x,y$ માટે $Rxy$, $x=y$ અથવા $Ryx$ હોય એવો આંશિક ક્રમ \emph{રેખીય ક્રમ} છે. સામાન્ય રીતે આ ગુણધર્મ ધરાવતો સંબંધ \emph{સંબદ્ધ} કહેવાય છે.

સંબંધો જોડોના ગણ હોવાથી ગણોની ક્રિયાઓ તેમના પર પણ કરી શકાય છે. દાખલા તરીકે સંબંધોનો યોગગણ અને છેદગણ લઈ શકાય છે. સંબંધોને સાંકળીને તેમનો \emph{સાપેક્ષ ગુણાકાર} $R\mid S$ પણ લઈ શકાય છે. $R$ને પોતાની સાથે કોઈ પણ ધન સંખ્યામાં વખત સાંકળવાથી મળતા બધા સંબંધોનો યોગગણ $R$નો \emph{પરંપરિત સંવર} $R^+$ છે.
